\section{Energy increase under a dyadic split}\label{sec:dyadic}

We apply the matrix estimate to a fixed admissible triple and compare
its energy with the average energy of its children. For this section,
the inputs are bounded, compactly supported, and constant on products
of intervals in \(\cD_m\), for some \(m\in\Z\).

\subsection{Matrices for a local integral}

For \(\mathbf I=(I_0,I_1,I_2)\in\cA_k\), \(k\ge m\), let \(H_v\)
have an orthonormal basis indexed by the intervals of \(\cD_m\)
inside \(I_v\). Call these intervals atoms. In the probability measure
\(\abs{I_v}^{-1}\dd x_v\), each atom has weight
\(w_v(i)=2^{m-k}\).
Define \(W_v:H_{v+1}\to H_v\) by
\begin{equation}\label{eq:weighted-edge}
 W_v(i,j)=\sqrt{w_v(i)w_{v+1}(j)}\,F_v(i,j),
\end{equation}
where \(F_v(i,j)\) is the constant value on the corresponding atom
pair. Use these maps in \eqref{eq:row-star}, and define
\begin{equation}\label{eq:local-energy}
 e_v(\mathbf I)=\Phi(T_v)=\Phi(R_vR_v^*),
 \qquad e(\mathbf I)=\sum_{v=0}^2 e_v(\mathbf I).
\end{equation}

If \(k\ge m+1\), let \(D_v\) be the diagonal matrix of the
\(h_{I_v}\)-values on the atoms, and put
\[
 \tau(\mathbf I)=\tr(D_0W_0D_1W_1D_2W_2).
\]
The diagonal matrices are symmetric contractions. In the expanded
trace, each vertex weight occurs as two square roots, so
\(\tau(\mathbf I)\) is exactly the normalized integral of the product
of the three inputs and three Haar signs. Therefore
\begin{equation}\label{eq:trace-integral}
 L_{\mathbf I}(F_0,F_1,F_2)=2^{2k}\tau(\mathbf I).
\end{equation}

\subsection{The four children}

If \(b_v\in\{0,1\}\) chooses the left or right half of \(I_v\), the
child index is \(2n_{I_v}+b_v\). The child triple is admissible
precisely when \(b_0\oplus b_1\oplus b_2=0\). Equivalently, with
\(s_v=(-1)^{b_v}\), its signs satisfy \(s_0s_1s_2=1\).
Write \(\operatorname{ch}(\mathbf I)\) for these four children.
Under uniform choice among them, each sign is uniform and every pair
of signs is independent. Figure~\ref{fig:children} displays the
conditional averages used below.

\begin{figure}[htbp]
\centering
\begin{tikzpicture}[font=\small,
 box/.style={draw=black!55,rounded corners,inner sep=9pt,align=center},
 >=Stealth]
 \node (parent) at (0,1.4) {an admissible parent \(\mathbf I\)};
 \node[box,fill=blue!5] (plus) at (-2.5,0)
 {\(s_0=+1\)\\[3pt]
 \((+,+,+)\qquad(+,-,-)\)};
 \node[box,fill=blue!5] (minus) at (2.5,0)
 {\(s_0=-1\)\\[3pt]
 \((-,+,-)\qquad(-,-,+)\)};
 \draw[->] (parent) -- (plus);
 \draw[->] (parent) -- (minus);
 \node[align=center] at (0,-1.25)
 {\(\E[s_1\mid s_0]=\E[s_2\mid s_0]=0\),\qquad
   \(s_1s_2=s_0\).};
\end{tikzpicture}
\caption{The four admissible child sign triples, grouped by the
first sign. Conditional on \(s_0\), the two possibilities in its box
have equal probability. Each neighboring sign is unbiased, but their
product is fixed. The sum of the two incident-edge Grams in
\eqref{eq:child-gram} uses only these marginal averages.}
\label{fig:children}
\end{figure}


\begin{proposition}[Local energy increase]\label{prop:increase}
For \(\mathbf I\in\cA_k\), \(k\ge m+1\), define
\[
 \delta(\mathbf I)
 =\frac14\sum_{\mathbf I'\in\operatorname{ch}(\mathbf I)}
                     e(\mathbf I')-e(\mathbf I).
\]
Then \(\delta(\mathbf I)\ge0\), and
\begin{equation}\label{eq:local-control}
 \abs{\tau(\mathbf I)}
       \le\frac{10\sqrt2}{3}\,\delta(\mathbf I).
\end{equation}
\end{proposition}

\begin{proof}
We keep the parent spaces fixed and extend each child matrix by zero.
Set
\[
 C_v=(\Id+s_vD_v)^{1/2}.
\]
It is \(\sqrt2\) times the projection onto the atoms in the chosen
half. Child probability weights double, so the extended child edge
matrix is \(C_vW_vC_{v+1}\). Its row star is
\[
 R_v(s)=C_vR_v\diag(C_{v+1},C_{v-1}).
\]
In particular,
\begin{equation}\label{eq:child-gram}
 R_v(s)R_v(s)^*
 =C_v\bigl(
 W_vC_{v+1}^2W_v^*
 +W_{v-1}^*C_{v-1}^2W_{v-1}
 \bigr)C_v.
\end{equation}
Fix \(s_v\). Each neighboring sign is still uniform, so its squared
multiplier has conditional expectation \(\Id\).
Although the neighbors are correlated, \eqref{eq:child-gram} has
no product between their multipliers. Thus
\[
 \E\bigl[R_v(s)R_v(s)^*\mid s_v\bigr]
       =C_vR_vR_v^*C_v.
\]
By the convexity in Lemma~\ref{lem:hessian},
\begin{equation}\label{eq:conditional}
 \E\bigl[e_v(\mathbf I(s))\mid s_v\bigr]
       \ge\Phi(C_vR_vR_v^*C_v)
       =\Phi(T_v+s_vU_v),
 \qquad U_v=R_v^*D_vR_v.
\end{equation}
The equality compares the two Gram matrices of \(C_vR_v\):
their nonzero eigenvalues agree, even though they act on different
spaces.

Since \(D_v\) is a symmetric contraction, \(-T_v\le U_v\le T_v\).
Averaging \eqref{eq:conditional} over \(s_v=\pm1\) and applying
Corollary~\ref{cor:gap} gives
\[
 \frac14\sum_{\mathbf I'\in\operatorname{ch}(\mathbf I)}
                     e_v(\mathbf I')-e_v(\mathbf I)
       \ge\frac{1}{2\sqrt2}b_{T_v}(U_v).
\]
Summing over \(v\) proves nonnegativity of \(\delta\).
Proposition~\ref{prop:trace} now gives
\[
 \abs{\tau(\mathbf I)}
 \le\frac53\sum_v b_{T_v}(U_v)
 \le\frac{10\sqrt2}{3}\delta(\mathbf I),
\]
as asserted.
\end{proof}
